\chapter{Diseñados para amar}

\epigraph{Dios es amor, y quien permanece en el amor permanece en Dios.}{1 Juan 4,16}

Hagamos una pausa.

Llevamos cinco capítulos caminando juntos, y quizá convenga detenerse un momento para mirar el paisaje que hemos recorrido. No para repetirlo, sino para verlo de golpe, como quien sube a una colina y contempla el camino que serpentea a sus pies.

\section{Lo que el cuerpo nos ha dicho}

Empezamos por el parto. Vimos que nacemos de la forma más difícil del reino animal, radicalmente inmaduros, con la cabeza sin cerrar y sin poder sostenernos solos. Que esa fragilidad no es un fallo sino el precio de nuestro cerebro extraordinario. Y que, precisamente por nacer así, necesitamos que alguien nos sostenga, nos alimente, nos quiera ---durante años--- para llegar a ser lo que somos.

Después descubrimos que el cuerpo no nos deja solos ante esa tarea. La oxitocina inunda a la madre en el parto y la lactancia, empujándola al vínculo. El padre cambia también: baja la testosterona, sube la oxitocina, se activan circuitos cerebrales que antes dormían. El cuerpo de ambos se transforma para cuidar.

Vimos que nos enamoramos con una química que se parece a la obsesión ---dopamina, noradrenalina, serotonina baja--- y que después, si el vínculo madura, la oxitocina y la vasopresina construyen algo más estable, más hondo: un apego que nos empuja a quedarnos. Que el amor no es solo un sentimiento bonito, sino que tiene un sustrato biológico que nos inclina a la permanencia.

Desmontamos el mito del amor romántico de las películas y descubrimos que el amor real no se encuentra: se construye. Que es decisión además de sentimiento. Que tiene capas ---deseo, amistad, entrega--- y que crece cuando se elige cada día, no cuando se espera que todo fluya solo.

Y vimos, por último, que la diferencia entre varón y mujer no es un accidente ni un problema, sino la condición misma del encuentro. Que solo puedo encontrarme de verdad con quien es distinto de mí. Que la alteridad no amenaza al amor: lo hace posible.

Todo esto lo hemos visto sin necesidad de abrir la Biblia, sin citar un solo dogma, sin invocar ninguna autoridad religiosa. Lo hemos leído en los huesos, en las hormonas, en las neuronas. En el cuerpo.

Y el cuerpo dice, con una claridad que asombra: estás hecho para la relación. Estás configurado para el vínculo. Estás inclinado al amor.

\section{La pregunta que faltaba}

Pero hay una pregunta que el cuerpo, por sí solo, no puede responder.

¿Por qué?

¿Por qué estamos hechos así? ¿Por qué la evolución ---o lo que sea que haya detrás--- nos ha configurado como seres que necesitan amar y ser amados para sobrevivir, para crecer, para ser plenamente humanos?

La biología ofrece una respuesta, y es honesta: porque así sobrevivimos mejor. Porque las crías cuidadas por padres vinculados tenían más probabilidades de llegar a la edad adulta. Porque la selección natural favoreció a los grupos que cooperaban, que cuidaban, que permanecían juntos. Es una respuesta verdadera. Pero ¿es suficiente?

Pensemos un momento. Si todo se reduce a supervivencia, entonces el amor es un truco ---un mecanismo elegante, sí, pero un truco al fin y al cabo---. La oxitocina que inunda a una madre cuando abraza a su hijo sería solo un recurso para que la especie no se extinga. El deseo de permanecer junto a alguien, de construir una vida compartida, de ser fiel, sería una ilusión útil. Una estrategia de la evolución disfrazada de sentimiento.

¿De verdad es eso todo?

Hay algo en nosotros que se rebela contra esa explicación. Algo que intuye que el amor que sentimos ---el de verdad, el que nos desborda, el que nos hace mejores, el que a veces duele--- es algo más que bioquímica al servicio de la reproducción. Que cuando una madre mira a su hijo recién nacido y siente que daría la vida por él, ahí hay algo que trasciende la selección natural. Que cuando dos personas se prometen fidelidad para siempre, están haciendo algo que ningún gen les ha pedido.

La ciencia puede describir el cómo. Puede explicar los mecanismos, las hormonas, los circuitos neuronales. Pero el porqué último ---la razón de que existamos así, configurados para el amor--- queda fuera de su alcance. No porque la ciencia falle, sino porque esa pregunta pertenece a otro orden.

Y es aquí donde este libro da un paso.

\section{El que nos diseña es Amor}

Hay una frase en la primera carta de san Juan que, cuando se lee despacio, lo cambia todo. Una frase breve, de apenas tres palabras en griego ---\textit{ho theós agápe estín}---, que contiene quizá la afirmación más radical de toda la historia del pensamiento:

\begin{quotation}
\textit{Dios es amor.}\footnote{1~Jn 4,8. El texto completo dice: ``El que no ama no ha conocido a Dios, porque Dios es amor.'' No dice que Dios \textit{tiene} amor, o que \textit{da} amor, o que \textit{manda} amar. Dice que Dios \textit{es} amor. Es una afirmación ontológica: define la naturaleza misma de Dios.}
\end{quotation}

No dice que Dios es poderoso (que lo es). No dice que Dios es sabio (que lo es). Dice que Dios es amor. No como una cualidad entre otras, sino como definición de lo que Él es. Amor no es algo que Dios hace; es lo que Dios es.

Y ahora viene la segunda pieza. El libro del Génesis dice que Dios creó al ser humano ``a su imagen y semejanza''.\footnote{Gn 1,27. La expresión hebrea \textit{tselem} (imagen) y \textit{demut} (semejanza) indica que el ser humano refleja algo de Dios mismo, no en un sentido físico, sino en su capacidad de conocer, de querer y de entregarse. Cf. Concilio Vaticano~II, Constitución pastoral \textit{Gaudium et Spes}, n.~12: ``Creyentes y no creyentes están generalmente de acuerdo en este punto: todos los bienes de la tierra deben ordenarse en función del hombre, centro y cima de todos ellos.''} Si Dios es amor, entonces su imagen ---el ser humano--- tiene que ser un ser capaz de amar. No puede ser de otro modo. Un espejo refleja lo que tiene delante; y si lo que tiene delante es amor, el reflejo será un ser hecho para amar.

Ahora volvamos la vista a todo lo que hemos recorrido en estos capítulos. La pelvis que se estrecha. El cerebro que crece. El bebé que nace indefenso. La oxitocina que inunda a la madre. La testosterona del padre que baja. La química que nos enamora. El apego que nos hace quedarnos. La diferencia que nos abre al otro.

¿Y si todo eso no fuera solo evolución?

¿Y si fuera la firma del Diseñador?

Quiero ser honesto: lo que acabo de hacer no es una demostración científica, sino una invitación. Los hechos biológicos que hemos recorrido juntos son sólidos y se sostienen por sí mismos, creas o no en un Diseñador. Pero hay un momento en que uno mira todo eso y siente que puede significar algo más ---y dar ese paso es un acto de confianza personal, no una conclusión de laboratorio---. Aquí es donde entra mi fe, y tú eres libre de acompañarme o de quedarte con la ciencia a solas, que ya dice mucho.

No digo que la evolución no sea real ---lo es, y lo hemos visto con rigor---. Digo que la evolución, siendo verdadera, puede no ser la explicación última. Que detrás del cómo biológico puede haber un porqué más profundo. Que las mismas fuerzas que la ciencia describe con precisión pueden ser, al mismo tiempo, el modo en que un Dios que es amor ha querido configurar a la criatura hecha a su imagen.

Estamos diseñados para amar. No solo porque la selección natural lo favoreció, sino porque Alguien que es Amor nos pensó así desde el principio.\footnote{El Catecismo de la Iglesia Católica lo expresa con claridad: ``Siendo a imagen de Dios, el ser humano tiene la dignidad de persona; no es solamente algo, sino alguien. Es capaz de conocerse, de poseerse y de darse libremente y entrar en comunión con otras personas'' (CIC, n.~357). Y el Concilio Vaticano~II: ``El hombre, única criatura terrestre a la que Dios ha amado por sí misma, no puede encontrar su propia plenitud si no es en la entrega sincera de sí mismo a los demás'' (\textit{Gaudium et Spes}, n.~24).}

\section{La libertad, condición del amor}

Pero Dios no se conformó con hacernos capaces de amar. Hizo algo más arriesgado, más hermoso y más terrible a la vez: nos hizo libres.

¿Por qué? Porque un amor obligado no es amor. Es obediencia, es rutina, es a lo sumo costumbre. Pero no es amor. El amor, para ser amor de verdad, necesita poder decir que no.

Pensemos en algo sencillo. Si alguien te dice ``te quiero'' porque no tiene más remedio, porque está programado para decirlo, porque no puede evitarlo... ¿te sentirías querido? Probablemente no. Sentirías que eres el destinatario de un automatismo, no de un afecto. Lo que hace que un ``te quiero'' valga es que la persona que lo dice podría no decirlo. Podría irse, podría callarse, podría elegir a otro. Pero te elige a ti. Libremente.

Dios quiere ser amado así. No por obligación, no por inercia, no por programación genética. Libremente. Y para eso nos da algo que ningún otro ser en la creación tiene de la misma manera: conciencia y libertad.\footnote{Cf. \textit{Gaudium et Spes}, n.~17: ``La dignidad humana requiere, por tanto, que el hombre actúe según su conciencia y libre elección, es decir, movido e inducido por convicción interna personal y no bajo la presión de un ciego impulso interior o de la mera coacción externa.''} La capacidad de conocer lo que hacemos y de elegir hacerlo ---o no hacerlo---.

¿Es esto un riesgo? Claro que lo es. La libertad implica que podemos elegir mal. Que podemos dañar, mentir, abandonar, traicionar. Que podemos mirar al amor de frente y darle la espalda. La historia humana está llena de libertad mal usada, y las heridas que deja son reales.

Pero la alternativa sería peor. La alternativa sería un mundo de autómatas que ``aman'' porque no pueden hacer otra cosa. Un mundo sin traiciones, sí, pero también sin fidelidad ---porque la fidelidad solo tiene sentido cuando podrías ser infiel---. Un mundo sin egoísmo, pero también sin generosidad ---porque la generosidad solo existe cuando podrías quedarte con todo---. Un mundo sin dolor, quizá, pero también sin alegría verdadera.

Dios no crea la libertad a regañadientes, como un defecto inevitable. La crea como condición necesaria. Sin libertad, no habría imagen de Dios. Sin libertad, no habría amor. Sin libertad, seríamos otra cosa ---quizá más eficiente, quizá más segura--- pero no seríamos humanos.\footnote{Santo Tomás de Aquino desarrolla esta idea con precisión: el ser humano es imagen de Dios precisamente en cuanto es principio de sus propias acciones, al tener libre albedrío y dominio sobre sus actos. Cf. \textit{Summa Theologiae}, I-II, prólogo.}

\section{La primera pregunta}

Y aquí es donde todo lo que hemos dicho aterriza en algo concreto. Tan concreto como una iglesia, dos personas de pie ante un altar, y un sacerdote que abre la boca para hacer la primera pregunta del rito del matrimonio.

No pregunta: ``¿Os queréis?'' No pregunta: ``¿Estáis enamorados?'' No pregunta: ``¿Prometéis ser fieles?'' Todo eso vendrá después. Lo primero, lo que va antes de cualquier otra cosa, es esto:

\begin{quotation}
\textit{¿Venís a contraer matrimonio sin ser coaccionados, libre y voluntariamente?}\footnote{Ritual del Matrimonio, n.~60 (edición típica en español). Las tres preguntas del escrutinio son, en este orden: libertad, fidelidad y apertura a la vida. El orden no es casual.}
\end{quotation}

¿Venís libremente?

La Iglesia lleva dos mil años celebrando matrimonios, y en todo ese tiempo no ha cambiado el orden de las preguntas. Primero la libertad. Antes que el amor, antes que la fidelidad, antes que los hijos. Porque sin libertad, nada de lo que sigue es real.

Si no vienes libremente, tu ``sí'' no es un sí. Es un sonido vacío. Si no vienes libremente, tu promesa de fidelidad es papel mojado. Si no vienes libremente, tu entrega es una representación teatral.

La libertad es el suelo sobre el que se construye todo lo demás. Y no es casualidad que sea así. Es coherencia perfecta con lo que somos: seres hechos a imagen de un Dios que es amor, y que por tanto solo pueden amar de verdad cuando eligen libremente hacerlo.

Pero conviene entender bien qué es la libertad, porque tendemos a confundirla con hacer lo que apetece en cada momento. Y no es eso. Pensad en alguien que decide correr una maratón. Desde el momento en que toma esa decisión, su vida cambia: se levanta temprano a entrenar, se somete a dietas, renuncia a salir, aguanta sesiones que no le apetecen nada. ¿Ha dejado de ser libre? Al contrario: es \emph{más} libre que antes, porque está usando su libertad para algo que ha elegido y que le trasciende. O pensad en el opositor que se encierra a estudiar durante dos, tres, cuatro años. Renuncia a vacaciones, a vida social, a dormir hasta tarde. ¿Es un esclavo? No. Es alguien que ha decidido libremente labrarse un futuro, y que acepta el precio porque la decisión es suya y el horizonte merece la pena.

Pues el matrimonio funciona exactamente igual. Elegir a una persona para toda la vida no te quita libertad: la pone al servicio de algo grande. Las renuncias que implica ---y las habrá, todos los días--- no son cadenas que te impone nadie. Son el entrenamiento del maratoniano, la disciplina del opositor: el precio libre de una decisión libre.

Decía san Josemaría Escrivá que ``porque quiero'' es la razón más sobrenatural que existe.\footnote{San Josemaría Escrivá de Balaguer (1902-1975), fundador del Opus Dei. Esta idea recorre su predicación sobre la libertad como fundamento de la vida cristiana. Cf. \textit{Amigos de Dios}, nn.~24-38.} Y tenía razón. Porque en esas dos palabras se manifiesta lo que nos hace imagen de Dios: la libertad. No ``porque debo'' ---que es obediencia, que es norma, que es el mínimo---. Sino ``porque quiero'' ---que es elección, que es don, que es lo más parecido a cómo actúa Dios---.

Y cuando dos personas se miran a los ojos en el altar y dicen ``sí, quiero'', están pronunciando la frase más sobrenatural que puede decir un ser humano. No ``sí, debo''. No ``sí, me toca''. ``Sí, quiero.'' Libremente. Como Dios nos ama: porque quiere, no porque necesite hacerlo.

\section{El puente}

Hemos llegado, entonces, a un lugar nuevo.

Durante la primera parte de este libro, hemos mirado al ser humano desde la ciencia: la biología, la neurociencia, la psicología, la antropología. Y hemos descubierto algo asombroso: que todo en nosotros ---huesos, hormonas, neuronas, emociones--- apunta en la misma dirección. Hacia el vínculo. Hacia el cuidado. Hacia el amor.

Ahora sabemos por qué. Estamos diseñados para amar porque el que nos diseña es Amor. Y nos hace libres para que ese amor sea verdadero.

Con esto en las manos, la segunda parte de este libro se abre de forma natural. Si estamos diseñados para amar ---si esa es nuestra vocación más profunda---, entonces la pregunta siguiente es inevitable: ¿cómo? ¿Cómo se concreta ese amor? ¿Qué forma toma? ¿Qué nos ayuda a vivirlo bien, y qué lo desfigura?

La fe cristiana tiene una propuesta. No la impone; la ofrece. Y lo que ofrece no es un conjunto de normas sobre lo que se puede y no se puede hacer, sino una mirada ---una forma de entender el amor, el matrimonio y la familia que, curiosamente, encaja con lo que el cuerpo ya nos había dicho desde el principio---.

Eso es lo que vamos a explorar ahora.

Un aviso antes de seguir. Si la primera parte de este libro ha podido sorprender ---porque quizá no esperabas encontrar oxitocina y topillos en un libro sobre matrimonio---, la segunda parte transitará un territorio más conocido: el de la fe cristiana, sus sacramentos, su visión del amor. Sobre esto hay mucho y muy bueno ya escrito ---de Juan Pablo II a Francisco, de C.\,S. Lewis a Benedicto XVI---. No pretendo inventar nada: pretendo reunir, con la misma honestidad de los capítulos anteriores, lo que otros han dicho mejor que yo, y ponerlo en diálogo con lo que el cuerpo ya nos había enseñado. Si la primera parte descubría, la segunda reconoce. Y reconocer algo que ya estaba ahí puede ser, a su manera, tan revelador como descubrirlo por primera vez.

\paracaminar

\begin{enumerate}
  \item Cuando leéis que ``Dios es amor'' y que estamos hechos a su imagen, ¿qué resuena en vosotros? ¿Os parece una idea lejana o conecta con algo que habéis experimentado?

  \item Hablad juntos sobre la libertad en vuestra relación. ¿Sentís que os elegís libremente cada día? ¿Hay algo ---presión social, inercia, miedo--- que a veces nubla esa libertad?

  \item La primera pregunta del rito matrimonial es ``¿Venís libremente?'' ¿Por qué creéis que va antes que las demás? ¿Qué os dice eso sobre lo que la Iglesia entiende por amor?
\end{enumerate}

\vinetacapitulo{ch06}
